% Part: methods
% Chapter: induction
% Section: induction-on-N

\documentclass[../../../include/open-logic-section]{subfiles}

\begin{document}

\olfileid{mth}{ind}{inN}

\olsection{پر~$\Nat$ استقرا}

په تر ټولو ساده بڼه کښې استقرا يو داسې تخنيک دے چې د ټولو
طبيعي عددونو دپاره د نتيجو د ثابتولو دپاره کارېږي۔ دا پر دې
حقيقت تکيه کوي چې له $0$ څخه په پيل کولو او بيا بيا د~$1$ په
زياتولو بالاخره هر طبيعي عدد ته رسېږو۔ نو د دې د ثابتولو
دپاره چې يوه خبره د هر عدد دپاره رښتيا ده، کولاے شو (1)~ثابته
کړو چې د $0$ دپاره رښتيا ده، او (2)~وښيو چې هر کله د يو
عدد~$n$ دپاره رښتيا وي، نو د راتلونکي عدد~$n+1$ دپاره هم
رښتيا ده۔ که «عدد~$n$ خاصيت $P$ لري» په $P(n)$ لنډ کړو، او
«عدد~$k$ خاصيت $P$ لري» په $P(k)$، او داسې نور، نو د $P(n)$
استقرايي ثبوت، د ټولو $n \in \Nat$ دپاره، دا برخې لري:
\begin{enumerate}
\item د $P(0)$ ثبوت، او
\item دا ثبوت چې د هر~$k$ دپاره، که $P(k)$ وي نو $P(k+1)$ هم وي۔
\end{enumerate}
د دې د پوره روښانولو دپاره فرض کړئ چې (1) او~(2) دواړه لرو۔
بيا (1) راته وايي چې $P(0)$ رښتيا دے۔ که~(2) هم ولرو، نو
په ځانګړي ډول پوهېږو چې که $P(0)$ وي نو $P(0+1)$، يعنې
$P(1)$ هم وي۔ دا له دې عام بيان څخه چې «د هر~$k$ دپاره، که
$P(k)$ وي نو $P(k+1)$ هم وي»، د~$0$ په اېښودلو، د~$k$ پر ځاے،
راوځي۔ نو د مودوس پوننس په قاعده~$P(1)$ لرو۔ له (2) څخه
بيا، اوس $1$ په اخيستلو، د~$n$ پر ځاے،، لرو چې که $P(1)$ وي
نو~$P(2)$ هم وي۔ ځکه چې همدا اوس مو~$P(1)$ ثابت کړ، د مقدم
د اثبات په قاعده~$P(2)$ لرو۔ او همداسې مخ په وړاندې ځو۔
د هر عدد~$n$ دپاره د دې کار له $n$~ځله کولو وروسته بالاخره
تر~$P(n)$ رسېږو۔ نو (1) او~(2) په ګډه~$P(n)$ ثابتوي، د هر $n
\in \Nat$ دپاره۔

يوه بېلګه ګورو۔ فرض کړئ غواړو معلومه کړو چې د نرد په $n$
دانو څو بېلابېلې مجموعې غورځولاے شو۔ که څه هم ښايي بې‌معنا
ښکاره شي، له $0$ دانو پيل کوو۔ که هېڅ دانه نۀ وي، نو يوازې
يوه ممکنه مجموعه «غورځولاے» شئ: هېڅ ټکے نشته، چې مجموعه ئې~$0$
ده۔ نو د غورځولو د بېلابېلو ممکنو مجموعو شمېر~$1$ دے۔ که
يوازې يوه دانه ولرئ، يعنې $n=1$، نو له $1$ څخه تر~$6$ پورې
شپږ ممکنه ارزښتونه دي۔ په دوو دانو له $2$ څخه تر $12$ پورې
هره مجموعه غورځولاے شو، چې $11$ امکانات دي۔ په دريو دانو
له $3$ څخه تر $18$ پورې هر عدد غورځولاے شو، يعنې $16$
بېلابېل امکانات۔ $1$، $6$، $11$، $16$: داسې ښکاري چې يو
نظم لري؛ ښايي ځواب $5n+1$ وي؟ البته $5n+1$ تر ټولو زيات
ممکن شمېر دے، ځکه چې يوازې $5n+1$ عددونه شته، د $n$، يعنې په $n$ دانو د غورځولو
تر ټولو ټيټ ارزښت چې ټولې $1$ وي، او د $6n$، يعنې تر ټولو
لوړ ارزښت چې ټولې $6$ وي، تر منځ۔

\begin{thm}
په $n$ نرد دانو ټول $5n+1$ ممکنه ارزښتونه غورځولاے شو،
د $n$ او $6n$ تر منځ۔
\end{thm}

\begin{proof}
$P(n)$ دې دا دعوه وي چې «د $n$ او~$6n$ تر منځ هر عدد
غورځول ممکن دي، په $n$~نرد دانو»۔ د استقرا د کارولو دپاره
دا ثابتوو:
\begin{enumerate}
\item \emph{د استقرا بنسټ} $P(1)$، يعنې په يوازې يوه دانه
  د $1$ او $6$ تر منځ هر عدد غورځولاے شئ۔
\item \emph{د استقرا ګام}، چې د ټولو $k$ دپاره، که $P(k)$
  وي نو~$P(k+1)$ هم وي۔
\end{enumerate}

(1) د نرد د يوې $6$-مخيزې دانې په کتلو ثابتېږي۔ دا ټول
6 مخونه لري، او د $1$ او~$6$ تر منځ هر عدد ئې پر يو مخ
ښکاري۔ نو په يوه دانه د $1$ او~$6$ تر منځ هر عدد غورځول
ممکن دي۔

د (2) د ثابتولو دپاره د شرط مقدم، يعنې $P(k)$، فرض کوو۔
دې فرض ته \emph{استقرايي فرض} وايي۔ له دې څخه د~$P(k+1)$
د ثابتولو دپاره کار اخلو۔ سخته برخه دا ده چې د $k+1$ دانو
د غورځولو د ممکنو ارزښتونو په هکله د فکر کولو داسې لاره
ومومو چې د $k$~دانو د غورځولو ممکن ارزښتونه او د اضافي
$k+1$-مې دانې د غورځولو ارزښتونه سره يوځاے کړي۔ خو که
استقرايي فرض کارول غواړو، نو همدا کار مو پکار دے۔

استقرايي فرض وايي چې د $k$ او~$6k$ تر منځ هر عدد ترلاسه
کولاے شو، په $k$~دانو۔ که~$1$ وغورځوو، په خپلې $(k+1)$-مې
دانې، نو مجموعې ته $1$ ورزياتوي۔ نو د $k+1$ او $6k+1$
تر منځ هر ارزښت غورځولاے شو، د $k$~دانو په غورځولو او
بيا د~$1$ په غورځولو، په $(k+1)$-مه دانه۔ څه پاتې شول؟ له $6k+2$
څخه تر $6k+6$ پورې ارزښتونه۔ دا په $k$ دانو د $6$ په
غورځولو او بيا د $2$ او $6$ تر منځ د يو عدد په غورځولو،
په خپلې $(k+1)$-مې دانې، ترلاسه کولاے شو۔ په ګډه دا
معنا لري چې په $k+1$ دانو د $k+1$ او $6(k+1)$ تر منځ
هر عدد غورځولاے شو۔ يعنې~$P(k+1)$ مو ثابت کړ، د~$P(k)$
فرض په کارولو، چې استقرايي فرض دے۔
\end{proof}

ډېر ځله استقرا هغه وخت کاروو چې د داسې څيزونو، لکه
عددونو، سټونو او نورو، د يوې لړۍ په هکله يوه خبره ثابتول
غواړو چې پخپله «په استقرايي ډول» تعريف شوې وي؛ يعنې
$(n+1)$-م څيز ئې د $n$-م له مخې تعريف شوے وي۔ د بېلګې
په توګه، د طبيعي عددونو مجموعه~$s_n$، تر~$n$ پورې، داسې
تعريفولاے شو:
\begin{align*}
  s_0 & = 0\\
  s_{n+1} & = s_n + (n+1)
\end{align*}
له دې تعريفه ترلاسه کېږي:
\begin{align*}
  s_0 & = 0,\\
  s_1 & = s_0 + 1 && = 1,\\
  s_2 & = s_1 + 2 && = 1 + 2 = 3\\
  s_3 & = s_2 + 3 && = 1 + 2 + 3 = 6, \text{ او داسې نور۔}
\end{align*}
اوس په استقرا ثابتولاے شو چې $s_n = n(n+1)/2$۔

\begin{prop}
  $s_n = n(n+1)/2$.
\end{prop}

\begin{proof}
بايد (1) دا ثابت کړو چې $s_0 = 0\cdot(0 + 1)/2$، او (2)
که $s_k =
  k(k+1)/2$ وي نو $s_{k+1} = (k+1)(k+2)/2$ هم وي۔ (1)
ښکاره دے۔ د (2) د ثابتولو دپاره استقرايي فرض اخلو چې
$s_k = k(k+1)/2$۔ د دې په کارولو بايد وښيو چې
$s_{k+1} = (k+1)(k+2)/2$۔

$s_{k+1}$ څه دے؟ د تعريف له مخې، $s_{k+1} = s_k + (k+1)$۔
د استقرايي فرض له مخې، $s_k = k(k+1)/2$۔ دا په تېره
معادله کښې ځايناستي کولاے شو، او بيا يوازې د کسرونو لږ
حساب ته اړتيا لرو:
  \begin{align*}
    s_{k+1} & = \frac{k(k+1)}{2} + (k+1) = {}\\
    & = \frac{k(k+1)}{2} + \frac{2(k+1)}{2} = {}\\
    & = \frac{k(k+1) + 2(k+1)}{2} = {}\\
    & = \frac{(k+2)(k+1)}{2}.
  \end{align*}
\end{proof}

دلته مهمه زده کړه دا ده چې که د يوې په استقرايي ډول
تعريف شوې لړۍ $a_n$ په هکله يوه خبره ثابتوئ، نو استقرا
ښکاره مناسبه لاره ده۔ او که لړۍ داسې تعريف شوې هم نۀ
وي، لکه د نرد د غورځولو د امکاناتو په بېلګه کښې، بيا
هم استقرا کارولاے شئ، که په څه ډول د~$k+1$ حالت د~$k$
له حالت سره تړلاے شئ۔

\end{document}
